% Chapter 3: Build Systems
% File: chapters/chapter-03-build/chapter-03-build.tex

\chapter{Build Systems and Cross-Compilation Setup}
\label{ch:chapter-03}
\chapthumb{3}

\begin{quote}
    \textit{``A build that succeeds is not evidence that it built what you wrote.''}
\end{quote}

Cross-compiling for a microcontroller means four things have to agree: the target triple, the memory
layout, the linker's idea of what a loadable segment is, and the image the programmer writes. When
they disagree, the build still succeeds — that is the nature of linker scripts — and the board either
does nothing or, in the worst case, no longer boots.

This chapter is the whole chain for this board, with the two build defects that produced a green log
and a wrong artefact.

\minitoc

\clearpage

\section{The Target Triple, Decoded}
\label{sec:build-target}

\codeinline{thumbv8m.main-none-eabihf} is four decisions in one string, and three of them can be
wrong in ways that still compile.

\begin{table}[htbp]
    \centering
    \footnotesize
    \begin{tabular}{@{}p{0.22\textwidth}p{0.30\textwidth}p{0.42\textwidth}@{}}
        \toprule
        \textbf{Part} & \textbf{Means} & \textbf{If you get it wrong} \\
        \midrule
        \codeinline{thumbv8m} & ARMv8-M | ARMv8-M Baseline & \codeinline{armv7m} links, then traps on an instruction the core has \\
        \codeinline{main} & mainline, not baseline & baseline drops the DSP and FPU instructions you may be using \\
        \codeinline{none} & no OS, no libc & \codeinline{std} becomes available and fails deep in linking \\
        \codeinline{eabihf} & hard-float ABI & \codeinline{eabi} links and quietly passes floats in the wrong registers \\
        \bottomrule
    \end{tabular}
    \caption{The four decisions. Only the last one fails \emph{quietly}.}
    \label{tab:target-triple}
\end{table}

\subsection{Hard Float, and Why It Matters Here}

The EFR32MG24 has a single-precision FPU (FPv5-SP). \codeinline{eabihf} tells the compiler it may
pass floating-point arguments in \codeinline{s0}--\codeinline{s15} and may emit FPU instructions
directly. With \codeinline{eabi} the same source compiles to calls into soft-float helpers, links
perfectly, and runs several times slower on exactly the arithmetic this project's classifier uses.

There is one trap on the way in: \codeinline{-mfpu} and \codeinline{-mfloat-abi} flags that
\emph{conflict} with the target's own choice produce a warning on every build — ``switch
\codeinline{-mcpu=cortex-m33} conflicts with switch \codeinline{-march=armv8-m.main+fp}''. That
warning is not harmless: it appeared once per C file in this project and buried the real warnings
under it. The flags belong to the target; a build script adding its own is duplicating a decision it
does not own.

\subsection{The Toolchain, in Files}

\begin{minted}{bash}
# the target must be installed once
rustup target add thumbv8m.main-none-eabihf

# .cargo/config.toml decides the default target and the runner
[build]
target = "thumbv8m.main-none-eabihf"

[target.thumbv8m.main-none-eabihf]
rustflags = ["-C", "link-arg=-Tlink.x"]
\end{minted}

\codeinline{link.x} comes from \codeinline{cortex-m-rt}, and it is the file that reads
\codeinline{memory.x} and lays the image out. That indirection is why a wrong \codeinline{memory.x}
produces no error: \codeinline{link.x} simply places the sections where it was told to.

\subsection{Per-Crate Features, Not Per-Crate Flags}

A workspace with several crates that must not duplicate the vector table uses a feature to decide
which \emph{one} of them pulls \codeinline{cortex-m-rt}:

\begin{minted}{toml}
# the HAL exposes the runtime but does not enable it by default
[features]
rt = ["dep:cortex-m-rt"]

# examples turn it on; there is exactly one #[entry] per binary
efr32mg24-hal = { path = "../efr32mg24-hal", features = ["rt"] }
\end{minted}

Without that split, two crates in one binary each define \codeinline{Reset} and the link fails with a
duplicate symbol that names neither crate. With it, the rule is explicit: a library may know about
the runtime; only a final binary may install it.

\section{The Bootloader Decides Your Memory Layout}
\label{sec:build-bootloader}

The single most dangerous line in an embedded project is the \codeinline{ORIGIN} of the flash region,
and it is usually copied from a tutorial for a board with no bootloader.

\subsection{What Is at the Base of Flash}

This board ships with a Gecko Bootloader occupying the first 24 KiB. On the EFR32MG24 flash begins at
\codeinline{0x08000000}, so:

\begin{table}[htbp]
    \centering
    \footnotesize
    \begin{tabular}{@{}p{0.26\textwidth}p{0.26\textwidth}p{0.42\textwidth}@{}}
        \toprule
        \textbf{Region} & \textbf{Addresses} & \textbf{Owner} \\
        \midrule
        Bootloader & \codeinline{0x08000000}--\codeinline{0x08005FFF} & factory; never written by this project \\
        Application & \codeinline{0x08006000}--end & every image in this book \\
        \bottomrule
    \end{tabular}
    \caption{24 KiB that a linker script must not touch.}
    \label{tab:flash-map}
\end{table}

\subsection{The Wrong memory.x, and What It Does}

The version that appeared in this project's first draft of the book — copied from a generic
\codeinline{cortex-m} example — read:

\begin{minted}{text}
MEMORY
{
  FLASH (rx) : ORIGIN = 0x00000000, LENGTH = 1536K
  RAM  (rwx) : ORIGIN = 0x20000000, LENGTH = 256K
}
\end{minted}

Every line is wrong for this board, and none of them produces an error. The linker happily lays the
image out from address zero, the programmer writes it exactly there, the write verifies — and the
board no longer boots, because the bootloader that was at zero has been replaced by an application
that expects to be at \codeinline{0x08006000}.

\subsection{The Version in Use}

\begin{minted}{text}
MEMORY
{
  /* The Gecko bootloader owns 0x08000000..0x08005FFF. An application that starts
     below 0x08006000 overwrites it. */
  FLASH (rx)  : ORIGIN = 0x08006000, LENGTH = 1536K - 24K
  /* The mailbox at 0x2003F000 is a convention shared by every example; the stack
     is placed below it so a deep call cannot scribble over the report. */
  RAM   (rwx) : ORIGIN = 0x20000000, LENGTH = 252K
}

_stack_start = ORIGIN(RAM) + LENGTH(RAM) - 4096;
\end{minted}

Both numbers are deliberate: 252K rather than 256K because the top 4 KiB are reserved, and
\codeinline{\_stack\_start} is below the mailbox so a stack overflow corrupts a report rather than the
data being reported.

\subsection{The Rule, Enforced in a Script}

A comment is not protection. The project has two mechanical guards:

\begin{minted}{bash}
# 1. where would this image land?
python3 scripts/elf-load-range.py target/thumbv8m.main-none-eabihf/release/examples/pac_led_blink
# load_min=0x08006000 load_max=0x08006523 size=1316 segments=3

# 2. refuse anything below the bootloader boundary, before writing
bash scripts/flash.sh "$ELF"          # dry run: prints the plan, refuses load_min < 0x08006000
bash scripts/flash.sh "$ELF" --confirm
\end{minted}

The first command is the one worth internalising: \textbf{read the load addresses before every flash},
not after a board stops answering. It takes one second and it is the difference between a mistake and
a brick.

\section{--nmagic, or the Linker's Idea of a Segment}
\label{sec:build-nmagic}

This section documents the least intuitive flag in the whole chain, because without it the image is
built, verified, and written to the wrong place.

\subsection{What the Linker Does by Default}

ELF loadable segments carry two addresses: a virtual address (\codeinline{VMA}, where the code expects
to run) and a physical address (\codeinline{LMA}, where the bytes must land in memory). For flash that
is always 1:1 — you cannot relocate code that runs from flash.

A linker aiming at an operating system still aligns segments to page boundaries, and to do that it
gives the segment an \emph{offset} that no longer matches its load address. An OS loader compensates;
a bare-metal programmer does not, and writes the file's bytes at the segment's file offset while the
code expects them at the segment's address.

\subsection{The Symptom, Measured}

\begin{minted}{text}
# with a default link: the loadable segment's file offset is 0x1000, its address is 0x08006000
probe-rs download image.elf        # writes bytes at the wrong offset
probe-rs verify image.elf          # verifies the same wrong offsets: it agrees with itself
                                   # the board does not run
\end{minted}

This is the failure that teaches the difference between ``verified'' and ``correct''. The programmer
compares the flash against the file it just wrote. If the file was written at the wrong offset by a
consistent rule, both sides agree and the check passes.

\subsection{The Fix, and Why It Is One Flag}

\begin{minted}{text}
# in .cargo/config.toml, or as a link argument
rustflags = ["-C", "link-arg=--nmagic"]
\end{minted}

\codeinline{--nmagic} tells the linker not to align segments and not to pad them — the image is flat,
VMA equals LMA for every loadable segment, and the programmer's naive write is correct. Every bare
metal Rust project needs it and most tutorials omit it, because their toolchain's runner happens to
compensate.

\subsection{Checking the Invariant Instead of Trusting the Flag}

The project's inspection script prints exactly what matters, so the flag is verified rather than
believed:

\begin{minted}{bash}
python3 scripts/elf-load-range.py "$ELF"
# for each PT_LOAD: paddr == vaddr, and load_min >= 0x08006000
#   seg paddr=0x08006158 vaddr=0x08006158 filesz=5400
#   seg paddr=0x08007670 vaddr=0x08007670 filesz=108
#   seg paddr=0x080076DC vaddr=0x20000000 filesz=4     <- the .data initialiser
\end{minted}

Three lines, and two facts read directly off them: every code segment has \codeinline{paddr == vaddr}
(the flag worked), and \codeinline{load\_min} is above the bootloader (the memory map is right). The
third segment is the initialiser for \codeinline{.data}, which correctly has a flash address and a RAM
address — that one is supposed to differ.

\section{Builds That Lie}
\label{sec:builds-that-lie}

Two defects in this project had the same shape, and it is the most dangerous shape a build system can
have: \textbf{the build succeeded, and the artefact was not what the sources said.}

\subsection{Case 1: A build.rs That Did Not Watch Its Own C File}

A C file compiled by the \codeinline{cc} crate, used to provide symbols the vendor blobs reference:

\begin{minted}{rust}
// build.rs — the version that lied
cc::Build::new().file("src/cmsis_stubs.c").compile("cmsis_stubs");
println!("cargo:rerun-if-changed=blobs/");
println!("cargo:rerun-if-changed=vendor/");
\end{minted}

Editing \codeinline{src/cmsis_stubs.c} changed nothing. Cargo had no reason to re-run the script, the
old object file stayed, and the link failed on a symbol that was visibly present in the source — with
a message that sends you looking at the blob instead of the build.

\begin{minted}{rust}
// the version that tells the truth
println!("cargo:rerun-if-changed=src/cmsis_stubs.c");
\end{minted}

The general rule: \textbf{any file a build script reads or compiles is an input, and an input it does
not declare is an input the build ignores.} Cargo does not scan your script for file paths.

\subsection{Case 2: A Book That Compiled Yesterday's Sources}

This book is built in a container on a different machine, from a tarball:

\begin{minted}{bash}
# the version that lied: extracts to efr32book/book/latex/
tar cf - book/latex | ssh host 'cd ~/efr32book && tar xf -'

# ...while the container compiles efr32book/latex/, so every build used the previous copy
podman exec --workdir /workspace/efr32book/latex tex xelatex main.tex
\end{minted}

The path \codeinline{book/latex} is preserved by \codeinline{tar}, the compile directory is not, and
the two never met. A whole section was written, ``built'' successfully, and did not appear in the PDF.

The fix is one flag plus one check:

\begin{minted}{bash}
tar cf - book/latex | ssh host 'cd ~/efr32book && tar xf - --strip-components=1'

# and then prove it, before compiling
md5sum book/latex/chapters/.../01.08-layer-boundaries.tex
ssh host 'md5sum ~/efr32book/latex/chapters/.../01.08-layer-boundaries.tex'   # must match
\end{minted}

\subsection{The Pattern Behind Both}

\begin{table}[htbp]
    \centering
    \footnotesize
    \begin{tabular}{@{}p{0.30\textwidth}p{0.28\textwidth}p{0.34\textwidth}@{}}
        \toprule
        \textbf{Lie} & \textbf{What you see} & \textbf{What catches it} \\
        \midrule
        unwatched input & link error naming a symbol you wrote & \codeinline{rerun-if-changed} on every input \\
        wrong extraction path & success, no new content & compare a hash before compiling \\
        cached artefact & old behaviour at new code & clean the package, or watch the log for the crate \\
        \bottomrule
    \end{tabular}
    \caption{Three ways to build the wrong thing, and the one-line check for each.}
    \label{tab:build-lies}
\end{table}

The common answer is the same in all three cases: \textbf{verify that the thing being built is the
thing you wrote.} It is a one-line check, and it is cheaper than any of the afternoons it saves.

\section{The Gate: What Runs Without a Board}
\label{sec:build-gate}

A project with hardware has a scheduling problem: the compiler is available constantly, the board is
available sometimes. Anything that needs the board gets deferred, and deferred checks do not happen.
The answer this project uses is a single script that runs everything that does not need hardware.

\subsection{What It Checks}

\begin{minted}{text}
$ bash scripts/gate.sh
PASS  cargo fmt --all --check
PASS  efr32mg24-hal clippy, 0 warnings
PASS  efr32mg24-pac clippy, 0 warnings
PASS  xiao-expansion-base-bsp clippy, 0 warnings
PASS  xiao-mg24-bsp clippy, 0 warnings
PASS  motion-detection clippy, 0 warnings
PASS  efr32mg24-hal: 5 examples of its own level
PASS  efr32mg24-pac: 19 examples of its own level
PASS  xiao-expansion-base-bsp: 9 examples of its own level
PASS  motion-detection: host tests (9 passed)
PASS  motion-detection: no hardware dependency
PASS  examples outside the PAC crate name no PAC type
PASS  no BSP re-exports the PAC
PASS  the HAL exposes no register block publicly
PASS  no example steals a core peripheral
VERDICT: PASS
\end{minted}

Every line is a property, not a command. That distinction matters: a script that runs
\codeinline{cargo clippy} tells you a tool ran; a line that says \emph{zero warnings in five crates}
tells you what is true.

\subsection{The Discovered-vs-Listed Rule}

The first version of this gate had a hard-coded list of crates, and the crate that most needed
checking was the one that was not on it. The version in use \emph{discovers} crates:

\begin{minted}{bash}
for dir in crates/*/ libraries/*/; do
  [ -f "${dir}Cargo.toml" ] || continue
  name=$(basename "$dir")
  ...
done
\end{minted}

A list is a snapshot of what existed when someone wrote it. Discovery cannot fall behind — which is
exactly what a check is for.

\subsection{The Rule Nobody Wrote Down}

The most valuable check in that output is the one that sounds odd: \textbf{every crate must have at
least one example of its own level}.

\begin{itemize}
    \item An API with zero examples passes every other check trivially — it compiles, it lints, and
          nothing calls it.
    \item It then rots silently: in this project a whole board module spent weeks documented and
          unbuildable, and the only symptom was that nobody used it.
    \item The check is cheap and it fails loudly, which is the entire point of a gate.
\end{itemize}

\subsection{What a Gate Cannot Do}

It cannot tell you the LED lit. This project's gate was green on a firmware that flashed successfully
and did not blink — four green steps, no light. That is not a failure of the gate; it is the boundary
of what a build-time check can know, and it is why the hardware checks in chapter~14 exist as a
separate, written procedure with its own measurements.


\clearpage

\begin{keytakeaways}
    \item The triple is \codeinline{thumbv8m.main-none-eabihf}: ARMv8-M mainline, no OS, hardware
          float. Getting the float ABI wrong links but produces a slow, subtly different binary.
    \item The bootloader owns the first 24 KiB. \codeinline{memory.x} starts flash at
          \codeinline{0x08006000}, and a script refuses any image that would start lower.
    \item \codeinline{--nmagic} is not optional: without it the linker emits a segment whose file
          offset disagrees with its load address, and the programmer writes to the wrong place.
    \item Two builds lied in this project: a \codeinline{build.rs} that did not watch its own C file,
          and a book build that compiled yesterday's sources because of a missing
          \codeinline{--strip-components}.
    \item A gate that runs without the board is worth more than a checklist, because it runs every
          time.
\end{keytakeaways}
